\section{Fixed-point structure and cycle decompositions}

\begin{theorem}[Positive definiteness under the right partial trace]
    \label{thm:posdef_partial_trace_right}
    \lean{Matrix.PosDef.partialTraceRight}
    \leanok
    Let $I$ and $J$ be finite index sets, and let
    $X\in\mathbb C^{(I\times J)\times(I\times J)}$ be positive definite.
    If $J$ is nonempty, then
    \[
        \bigl(\tr_{J}X\bigr)_{i,i'}
          =\sum_{j\in J}X_{(i,j),(i',j)}
    \]
    is positive definite.
\end{theorem}

\begin{proof}
    \leanok
    The right partial trace is the sum, over the nonempty set $J$, of the
    positive-definite principal submatrices of $X$ indexed by $I\times\{j\}$.
    A nonempty finite sum of positive-definite matrices is positive definite.
\end{proof}

\begin{theorem}[Positive definiteness under the left partial trace]
    \label{thm:posdef_partial_trace_left}
    \lean{Matrix.PosDef.partialTraceLeft}
    \leanok
    Let $I$ and $J$ be finite index sets, and let
    $X\in\mathbb C^{(I\times J)\times(I\times J)}$ be positive definite.
    If $I$ is nonempty, then
    \[
        \bigl(\tr_{I}X\bigr)_{j,j'}
          =\sum_{i\in I}X_{(i,j),(i,j')}
    \]
    is positive definite.
\end{theorem}

\begin{proof}
    \leanok
    The left partial trace is the sum, over the nonempty set $I$, of the
    positive-definite principal submatrices of $X$ indexed by $\{i\}\times J$.
    A nonempty finite sum of positive-definite matrices is positive definite.
\end{proof}

\begin{theorem}[Positive Kronecker factor from a normalized product]
    \label{thm:posdef_left_kronecker_trace_one}
    \lean{Matrix.PosDef.left_of_kronecker_of_trace_eq_one}
    \leanok
    Let $I$ and $J$ be nonempty finite index sets, with
    $A\in\mathbb C^{I\times I}$ and $B\in\mathbb C^{J\times J}$.  If
    \[
        A\otimes B>0
        \qquad\text{and}\qquad
        \tr(A)=1,
    \]
    then $A>0$.
\end{theorem}

\begin{proof}
    \leanok
    \uses{thm:posdef_partial_trace_left,
        thm:posdef_partial_trace_right}
    By Theorem~\ref{thm:posdef_partial_trace_left},
    \[
        \tr_{I}(A\otimes B)=\tr(A)B=B>0.
    \]
    Hence $\tr(B)>0$.  Theorem~\ref{thm:posdef_partial_trace_right} gives
    \[
        \tr_{J}(A\otimes B)=\tr(B)A>0.
    \]
    Rescaling by the positive number $\tr(B)^{-1}$ proves $A>0$.
\end{proof}

% Completed support transport recorded in
% docs/paper-gaps/wolf_theorem6_14_fixed_point_projection_gap.tex.
\begin{theorem}[Density blocks after restriction to maximal stationary support]%
    \label{thm:maximal_support_compression_density_blocks}
    \lean{IsPositiveMap.exists_block_densities_of_maximalSupportCompression,
        Matrix.traceAdjointMap_stationarySupportCompression,
        IsSchwarzMap.stationarySupportCompression,
        IsSchwarzMap.traceAdjointMap_stationarySupportCompression}
    \leanok
    \uses{thm:stationarySupport_compression,
        thm:mean_ergodic_projection_density_block_form}
    Let $T:\MN{D}\to\MN{D}$ be positive and trace preserving, and suppose
    that $T^*$ satisfies the Schwarz inequality.  Set
    $\rho_0=T_\infty(\Id)$ and let $V:\mathcal H\to\mathbb C^D$ be an
    isometry onto the support of $\rho_0$.  Then
    \[
        \widetilde T(Y)=V^*T(VYV^*)V
    \]
    has a positive definite fixed point.  Moreover, $T(B)=B$ if and only if
    there is a fixed point $Y$ of $\widetilde T$ such that $B=VYV^*$.
    There are positive integers
    $d_k,m_k$, a unitary $U$ on $\mathcal H$, and positive definite
    density matrices $\sigma_k\in\MN{m_k}$ such that
    \[
        \Fix(\widetilde T)
        =U\left(\bigoplus_k \sigma_k\otimes\MN{d_k}\right)U^*.
    \]
    This is the maximal-support application of the full-support step in the
    proof of \cite[Theorem~6.14]{Wolf2012Quantum}.  This intermediate theorem
    does not itself transport the displayed decomposition back to the original
    space and adjoin the complementary zero summand; that final step is
    Theorem~\ref{thm:wolf_6_14}.
\end{theorem}

\begin{proof}
    \leanok
    \uses{thm:stationarySupport_compression,
        thm:mean_ergodic_projection_density_block_form,
        thm:posdef_left_kronecker_trace_one}
    Theorem~\ref{thm:stationarySupport_compression} gives positivity, trace
    preservation, a positive definite fixed point for $\widetilde T$, and the
    stated correspondence between the original and compressed fixed points.
    Directly from the trace pairing,
    \[
        \widetilde T^*(A)=V^*T^*(VAV^*)V.
    \]
    Compressing the Schwarz defect for $T^*$ and adding the positive term
    through $\Id-VV^*$ proves the Schwarz inequality for
    $\widetilde T^*$.  Theorem~\ref{thm:mean_ergodic_projection_density_block_form}
    gives the fixed-point description with positive semidefinite trace-one
    matrices $\sigma_k$.

    Apply this description to a positive definite fixed point
    $\rho_{\mathrm{full}}$ of $\widetilde T$.  The unitary coordinate matrix
    is positive definite and has the form
    \[
        U^*\rho_{\mathrm{full}}U
          =\bigoplus_k\sigma_k\otimes X_k.
    \]
    Hence every principal block $\sigma_k\otimes X_k$ is positive definite.
    Both tensor factors have positive dimension, and $\tr(\sigma_k)=1$.
    Theorem~\ref{thm:posdef_left_kronecker_trace_one}, applied to
    $\sigma_k\otimes X_k$, therefore gives $\sigma_k>0$.
\end{proof}

\begin{theorem}[Fixed points of trace-preserving maps]%
    \label{thm:wolf_6_14}
    \lean{IsPositiveMap.exists_fixedPoints_densityBlocks_with_zero}
    \leanok
    \uses{thm:maximal_support_compression_density_blocks}
    Let $T:\MN{D}\to\MN{D}$ be positive and trace preserving, and suppose
    that $T^*$ satisfies the Schwarz inequality.  There are positive integers
    $d_k,m_k$, positive definite density matrices
    $\sigma_k\in\MN{m_k}$, a nonnegative integer $d_0$, and a unitary $U$
    associated with a decomposition
    \[
        \mathbb C^D=\mathbb C^{d_0}\oplus
          \bigoplus_k\mathbb C^{m_k}\otimes\mathbb C^{d_k}
    \]
    such that
    \[
        \Fix(T)
          =U\left(0_{d_0}\oplus
            \bigoplus_k\sigma_k\otimes\MN{d_k}\right)U^*.
    \]
    This is \cite[Theorem~6.14 and Equation~(6.63)]{Wolf2012Quantum}.
    The two tensor factors are interchanged from the displayed convention in
    the reference; the interchange is a unitary change of basis on each
    summand.
\end{theorem}

\begin{proof}
    \leanok
    \uses{thm:maximal_support_compression_density_blocks,
        thm:isometric_inclusion_zero_extension}
    Let $V:\mathbb C^n\to\mathbb C^D$ be the maximal-support isometry and
    let $U_c$ be the unitary in the density-block description of the
    compressed fixed points.  The fixed-space equivalence gives
    \[
        \Fix(T)
          =VU_c\left(\bigoplus_k
              \sigma_k\otimes\MN{d_k}\right)U_c^*V^*.
    \]
    Put $W=VU_c$.  Then $W^*W=\Id_n$.  Compare $W$ with the canonical
    inclusion of $\mathbb C^n$ as the second summand of
    $\mathbb C^{D-n}\oplus\mathbb C^n$.  The two inclusions have the same
    Gram matrix, so the finite-dimensional unitary extension theorem gives
    an ambient unitary $U$ carrying the canonical inclusion to $W$.  Hence,
    for every $A\in\MN{n}$,
    \[
        WAW^*=U(0_{D-n}\oplus A)U^*.
    \]
    Substituting the compressed density-block description for $A$ proves the
    assertion.  The first block is the whole orthogonal complement of the
    maximal stationary support, and is therefore one zero summand.  Positive
    definiteness and trace one for every $\sigma_k$ are supplied by
    Theorem~\ref{thm:maximal_support_compression_density_blocks}.
\end{proof}

\begin{theorem}[Conjugation by the square root at a fixed point of maximal support]%
    \label{thm:maximalSupport_weightedCorner}
    \lean{Kraus.exists_maximalSupport_weightedCorner_sqrt_eq}
    \leanok
    \uses{def:tp_kraus, def:support_proj}
    Let $T(X) = \sum_i K_i X K_i^\dagger$ be trace-preserving. There is a positive
    semidefinite fixed point $\rho_0$ of $T$ such that every fixed point $X$ of $T$
    arises as
    \[
        X \;=\; \sqrt{\rho_0}\, Y \sqrt{\rho_0}
    \]
    for a corner-supported $Y$ with $\sqrt{\rho_0}\, Y \sqrt{\rho_0}$ fixed by $T$.
\end{theorem}

\begin{proof}
    \leanok
    \uses{thm:maximalSupport_fixedPoint, thm:weightedCorner_sqrt_surjective}
    Take the fixed point $\rho_0$ of Theorem~\ref{thm:maximalSupport_fixedPoint}. Every
    fixed point $X$ of $T$ satisfies $Q_0\, X\, Q_0 = X$ for the support projection
    $Q_0$ of $\rho_0$, so Theorem~\ref{thm:weightedCorner_sqrt_surjective} applies to
    $X$ and produces the corner-supported $Y$.
\end{proof}

At the fixed point $\rho_0$, with the inverse square root taken on the support of
$\rho_0$, the carrier of the $*$-subalgebra of
Theorem~\ref{thm:weightedCornerFixedPointsStarSubalgebra} therefore realizes
\[
    \rho_0^{-1/2}\,\{X \in \MN{D} \mid T(X) = X\}\,\rho_0^{-1/2}:
\]
the conjugated
fixed-point set of Corollary~6.7 of~\cite{Wolf2012Quantum} at a fixed point of maximal
support, without a support restriction on the fixed points.

The corollary in~\cite{Wolf2012Quantum} quantifies over an arbitrary maximum-rank
fixed-point density matrix. The transfer from the constructed witness to every fixed
point of maximum rank goes through the rank of the support projection.

\begin{lemma}[The support projection has the rank of the matrix]%
    \label{lem:rank_stationaryProj}
    \lean{Kraus.rank_stationaryProj}
    \leanok
    \uses{def:support_proj}
    For $\rho \succeq 0$ on $\C^{D}$, the support projection of $\rho$ has the rank of
    $\rho$.
\end{lemma}

\begin{proof}
    \leanok
    \uses{def:support_proj}
    In the spectral decomposition $\rho = U\,\mathrm{diag}(\lambda)\,U^{\dagger}$ the
    support projection is $U\,\mathrm{diag}(\mathbf 1_{\lambda > 0})\,U^{\dagger}$.
    Multiplication by the invertible $U$ and $U^{\dagger}$ preserves the rank, and
    \[
        \rank\,\mathrm{diag}(\mathbf 1_{\lambda > 0})
        \;=\; \#\{i \mid \lambda_i > 0\}
        \;=\; \#\{i \mid \lambda_i \neq 0\}
        \;=\; \rank\,\mathrm{diag}(\lambda),
    \]
    the middle equality because the eigenvalues of $\rho \succeq 0$ are nonnegative.
\end{proof}

\begin{lemma}[The trace of the support projection is the rank]%
    \label{lem:trace_stationaryProj}
    \lean{Kraus.trace_stationaryProj}
    \leanok
    \uses{def:support_proj}
    For $\rho \succeq 0$ on $\C^{D}$, the trace of the support projection of $\rho$ is
    the rank of $\rho$.
\end{lemma}

\begin{proof}
    \leanok
    \uses{def:support_proj}
    With $\rho = U\,\mathrm{diag}(\lambda)\,U^{\dagger}$ as above,
    \[
        \tr\bigl(U\,\mathrm{diag}(\mathbf 1_{\lambda > 0})\,U^{\dagger}\bigr)
        \;=\; \tr\,\mathrm{diag}(\mathbf 1_{\lambda > 0})
        \;=\; \#\{i \mid \lambda_i > 0\}
        \;=\; \rank \rho,
    \]
    the last equality because the rank of a positive semidefinite matrix is the number
    of its positive eigenvalues.
\end{proof}

\begin{theorem}[The support of a fixed point of maximum rank carries every fixed
    point]%
    \label{thm:maximalSupport_of_maximalRank}
    \lean{Kraus.maximalSupport_of_maximalRank}
    \leanok
    \uses{def:tp_kraus, def:support_proj}
    Let $T(X) = \sum_i K_i X K_i^\dagger$ be trace-preserving and let $\rho \succeq 0$
    be a fixed point of $T$ whose rank bounds the rank of every fixed-point density
    matrix: $\rank \sigma \le \rank \rho$ for every
    $\sigma \succeq 0$ with $\tr \sigma = 1$ and $T(\sigma) = \sigma$. Then the support
    projection $Q$ of $\rho$ satisfies
    \[
        Q\, X\, Q = X
    \]
    for every fixed point $X$ of $T$.
    % Maintainer note: the maximum-rank fixed point of Corollary 6.7 of Wolf (2012) is a
    % density matrix; unit trace of the fixed point itself is not needed for the
    % conclusion and is not assumed here.
\end{theorem}

\begin{proof}
    \leanok
    \uses{thm:maximalSupport_fixedPoint, lem:rank_stationaryProj,
        lem:trace_stationaryProj}
    Let $\rho_0$ be the fixed point of maximal support of
    Theorem~\ref{thm:maximalSupport_fixedPoint}, with support projection $Q_0$, and let
    $P$ be the support projection of $\rho$. If $\rho_0 = 0$, then every fixed point
    vanishes, because $Q_0 = 0$ annihilates it, and the claim is trivial; so let
    $\rho_0 \neq 0$. Its trace is then a nonzero nonnegative real number, and
    $\sigma = (\tr \rho_0)^{-1} \rho_0$ is a fixed-point density matrix with
    $\rank \sigma = \rank \rho_0$, so the rank hypothesis
    gives $\rank \rho_0 \le \rank \rho$. Conversely
    $Q_0 \rho Q_0 = \rho$, so
    \[
        \rank \rho
        \;=\; \rank\bigl(Q_0\,\rho\,Q_0\bigr)
        \;\le\; \rank Q_0
        \;=\; \rank \rho_0,
    \]
    the last equality by Lemma~\ref{lem:rank_stationaryProj}; hence the ranks agree.
    From $\rho\,Q_0 = \rho$ one has $\rho\,(\Id - Q_0) = 0$, so the range of
    $\Id - Q_0$ lies in the kernel of $\rho$; the support projection $P$ of $\rho$
    vanishes on that kernel, so
    \[
        P\,(\Id - Q_0) = 0,
        \qquad P\,Q_0 = P,
    \]
    and taking conjugate transposes also $Q_0 P = P$. The difference $R = Q_0 - P$ is then Hermitian and idempotent,
    \[
        R^2 \;=\; Q_0 - P - P + P \;=\; R,
    \]
    and its trace vanishes by Lemma~\ref{lem:trace_stationaryProj} and the equality of
    the ranks, so
    \[
        \tr\bigl(R^{\dagger} R\bigr) \;=\; \tr R \;=\; 0,
    \]
    hence $R = 0$ and $P = Q_0$. The claim now follows from the maximal-support property
    of $\rho_0$.
\end{proof}

\begin{theorem}[Conjugation by the square root at every fixed point of maximum rank]%
    \label{thm:weightedCorner_sqrt_eq_of_maximalRank}
    \lean{Kraus.exists_weightedCorner_sqrt_eq_of_maximalRank}
    \leanok
    \uses{def:tp_kraus, def:support_proj}
    Let $T(X) = \sum_i K_i X K_i^\dagger$ be trace-preserving and let $\rho \succeq 0$
    be a fixed point of $T$ whose rank bounds the rank of every fixed-point density
    matrix. Then every fixed point $X$ of $T$ arises as
    \[
        X \;=\; \sqrt{\rho}\, Y \sqrt{\rho}
    \]
    for a corner-supported $Y$ with $\sqrt{\rho}\, Y \sqrt{\rho}$ fixed by $T$.
\end{theorem}

\begin{proof}
    \leanok
    \uses{thm:maximalSupport_of_maximalRank, thm:weightedCorner_sqrt_surjective}
    By Theorem~\ref{thm:maximalSupport_of_maximalRank} every fixed point $X$ of $T$
    satisfies $Q\, X\, Q = X$ for the support projection $Q$ of $\rho$, so
    Theorem~\ref{thm:weightedCorner_sqrt_surjective} applies to $X$ and produces the
    corner-supported $Y$.
\end{proof}

At every fixed point $\rho$ of maximum rank, with the inverse square root taken on the
support of $\rho$, the carrier of the $*$-subalgebra of
Theorem~\ref{thm:weightedCornerFixedPointsStarSubalgebra} therefore realizes the
conjugated fixed-point set
$\rho^{-1/2}\,\{X \in \MN{D} \mid T(X) = X\}\,\rho^{-1/2}$ of Corollary~6.7
of~\cite{Wolf2012Quantum}.

\section{Further irreducibility and primitivity equivalences}

This section collects several criteria and equivalences from~\cite[Theorems~6.2, 6.7, 6.8]{Wolf2012Quantum} that are
used later.

\begin{theorem}[Exponential positivity condition]%
    \label{thm:wolf_6_2_item_3}
    \lean{exp_posDef_of_irreducible_cp}
    \lean{irreducible_iff_exp_posDef_forall}
    \leanok
    \uses{def:irreducible_cp}
    Let $E$ be an irreducible completely positive map, $A \ge 0$ with
    $A \neq 0$, and $t > 0$. Then $\exp(tE)(A)$ is positive definite.
    Conversely, if $\exp(tE)(A)$ is positive definite for every such $t$, $A$,
    then $E$ is irreducible (full equivalence \lean{irreducible_iff_exp_posDef_forall}).
\end{theorem}

\begin{theorem}[Kraus-span equivalence]%
    \label{thm:wolf_6_8_kraus_span}
    \lean{MPSTensor.wolf_theorem_6_8_kraus_span}
    \leanok
    \uses{def:primitive_channel}
    Under normalization, primitivity in Wolf's sense is equivalent to eventual
    full Kraus rank.
\end{theorem}

\begin{theorem}[Combined primitivity equivalence]%
    \label{thm:wolf_6_8_conjunction}
    \lean{MPSTensor.wolf_theorem_6_8_conjunction}
    \leanok
    \uses{thm:wolf_6_8_kraus_span}
    Under normalization, primitivity in Wolf's sense is equivalent to the conjunction
    of eventual full Kraus rank, normality, and strong irreducibility.
\end{theorem}

\begin{remark}[Pairwise primitivity equivalences]
    \label{rem:wolf_6_8_pairwise}
    \lean{MPSTensor.primitivePaper_iff_hasEventuallyFullKrausRank}
    \lean{MPSTensor.primitivePaper_iff_stronglyIrreducible}
    \lean{MPSTensor.hasEventuallyFullKrausRank_iff_stronglyIrreducible}
    \lean{MPSTensor.hasEventuallyFullKrausRank_iff_isNormal}
    \leanok
    These are the pairwise equivalences in
    \cite[Theorem~6.8]{Wolf2012Quantum}.
\end{remark}

\begin{theorem}[Scalar fixed-point algebra case]
    \label{thm:wolf_6_15_scalar_conditional_expectation}
    \lean{Kraus.scalarConditionalExpectation_isConditionalExpectation}
    \leanok
    \uses{def:channel}
    In the scalar fixed-point algebra setting, the weighted map
    $X \mapsto \frac{\tr(\rho X)}{\tr(\rho)}\Id$ is a conditional
    expectation onto the adjoint fixed-point $*$-subalgebra.
\end{theorem}

\section{Multi-cycle block-permutation structure}

This section develops the block-permutation structure underlying the
cycle decomposition of~\cite[Theorem~6.16]{Wolf2012Quantum}. The
multi-cycle decomposition refines the single-cycle case to handle
arbitrary peripheral periods.

\begin{definition}[Block-permutation structure]
    \label{def:cycle_structure}
    \lean{CycleStructure}
    \leanok
    Let $T : \MN{D} \to \MN{D}$. A \emph{block-permutation structure} for
    $T$ consists of a finite index set $\iota$, a family of orthogonal
    projections $P_k \in \MN{D}$ for $k \in \iota$, and a permutation
    $\sigma \in \mathrm{Sym}(\iota)$ (not necessarily a single cycle---in
    general $\sigma$ may have multiple disjoint cycles), such that
    \[
        T(P_{\sigma(k)}) = P_k,\quad
        T(P_k X) = T(P_k)\,T(X),\quad
        T(X P_k) = T(X)\,T(P_k)
    \]
    for every $k \in \iota$ and $X \in \MN{D}$.
\end{definition}

\begin{definition}[Multi-cycle block-permutation structure]
    \label{def:multi_cycle_decomposition}
    \lean{MultiCycleDecomposition}
    \leanok
    A \emph{multi-cycle decomposition} of $T : \MN{D} \to \MN{D}$ consists of a
    finite cycle index $\iota$, a per-cycle period
    $m : \iota \to \N_{>0}$, a projection family
    $P_{c,k} \in \MN{D}$ for $c \in \iota$ and $k \in \Fin{m(c)}$
    consisting of orthogonal projections. For every $c \in \iota$ and
    $k \in \Fin{m(c)}$, the per-cycle cyclic action is
    \[
        T(P_{c,k+1}) = P_{c,k}.
    \]
    The multiplicative-domain factorisations are
    $T(P_{c,k}\,X) = T(P_{c,k})\,T(X)$ and
    $T(X\,P_{c,k}) = T(X)\,T(P_{c,k})$.
\end{definition}

\begin{lemma}[Per-cycle corner preservation]
    \label{thm:multi_cycle_preserves_corner_pow_period}
    \lean{MultiCycleDecomposition.preserves_corner_pow_period}
    \leanok
    \uses{def:multi_cycle_decomposition}
    For every cycle $c \in \iota$ and index $k \in \Fin{m(c)}$, the
    iterate $T^{m(c)}$ preserves the corner $P_{c,k} \MN{D}\,P_{c,k}$.
\end{lemma}

\begin{proof}
    \leanok
    \uses{def:multi_cycle_decomposition}
    By induction on $n$, the $n$-th iterate $T^n$ sends
    $P_{c,k+n}\,X\,P_{c,k+n}$ to $P_{c,k}\,T^n(X)\,P_{c,k}$, using the
    multiplicative-domain factorisations and the cyclic action
    $T(P_{c,j+1}) = P_{c,j}$ at each induction step. Specialising to
    $n = m(c)$ and using that the indices $k + m(c) = k$ in $\Fin{m(c)}$
    yields corner preservation of $T^{m(c)}$.
\end{proof}

\begin{lemma}[Common-period corner preservation]
    \label{thm:multi_cycle_preserves_corner_pow_of_dvd}
    \lean{MultiCycleDecomposition.preserves_corner_pow_of_dvd}
    \leanok
    \uses{thm:multi_cycle_preserves_corner_pow_period}
    If $N \in \N$ is divisible by every per-cycle period $m(c)$,
    then $T^N$ preserves every corner $P_{c,k} \MN{D}\,P_{c,k}$ of the
    decomposition.
\end{lemma}

\begin{proof}
    \leanok
    \uses{thm:multi_cycle_preserves_corner_pow_period}
    Write $N = m(c)\,q$. Corner preservation is stable under taking
    powers, so $T^N = (T^{m(c)})^q$ preserves each corner by
    Lemma~\ref{thm:multi_cycle_preserves_corner_pow_period}.
\end{proof}

\begin{definition}[Flattening to a single-index block-permutation structure]
    \label{def:multi_cycle_to_cycle_structure}
    \lean{MultiCycleDecomposition.toCycleStructure}
    \leanok
    \uses{def:multi_cycle_decomposition, def:cycle_structure}
    A multi-cycle decomposition gives a single-index block-permutation
    structure on the sigma index $\Sigma_{c \in \iota} \Fin{m(c)}$: the
    permutation is the sigma-product of per-cycle cyclic shifts
    $k \mapsto k+1$ (whose cycle decomposition has one cycle per
    $c \in \iota$), the projections are $(c, k) \mapsto P_{c,k}$, and the
    multiplicative-domain factorisations descend componentwise. The
    resulting map forgets the explicit cycle-indexing.
\end{definition}
